\chapter{Basic Commands}
\section{Defining the Basic Geometry of Calculation}
To define the basic geometry type of calculation, the command
\begin{verbatim} 
<geometry type="plane">
\end{verbatim} 
is used. Allowed parameters for {\bf type} are
{\bf 
\begin{itemize}
\item plane
\item axi
\item 3d
\end{itemize}}

If needed, a more detailed description of the geometry type must be
provided in the individual PDEs (e.g. in mechanics: {\bf "planeStrain"}).



\section{Defining the Input and Output File Format}
The format of the input and ouput files must be defined by the
commands {\bf input} and {\bf output}:

\begin{verbatim} 
<input  format="mesh" meshLibrary="cfsGrid"/>
<output format="gmv"/>
\end{verbatim} 

Allowed values of {\bf format} for input files are
{\bf \begin{itemize}
\item mesh
\item dat
\end{itemize}}
and for output files
{\bf \begin{itemize}
\item unv
\item gmv
\end{itemize}}
while the meshLibrary attribute can currently be one of
\begin{itemize}
\item \textbf{cfsGrid }
\item \textbf{adaptGrid}
\end{itemize}
\noindent
By default, the input format is \textbf{mesh}, the ouput format is \textbf{unv}
and the value for meshLibrary is \textbf{cfsGrid }.


\section{Defining the Material File}
To specify the name of the material file, the command 
\begin{verbatim}
<materialData file="matFileName"/>
\end{verbatim}
must be used. The default name for the material file is {\bf mat.dat}.


\section{Analysis Techniques}\label{subsec:Analysis}
In the config-file, the command defining the analysis type reads as

{\bf $<$analysis type=''TYPE''/$>$}

\noindent
where {\bf TYPE} can be chosen one of
{\bf
\begin{itemize} 
\item transient 
\item static
\item harmonic
\item eigenfrequency
\end{itemize}}


\subsection{Static Analysis}
For static analysis, no additional information beside the command

{\bf \verb!<analysis type="static"/>!}

\noindent
are needed.

\subsection{Transient Analysis}

To be able to define the transient analysis in more detail, several
additional parameters have to be used:

\begin{verbatim}
<transient>
   <numSteps>    20   </numSteps>
   <firstDt>     5e-5 </firstDt>
   <stepSaveBeg> 1    </stepSaveBeg>
   <stepSaveEnd> 20   </stepSaveEnd>
   <stepSaveInc> 1    </stepSaveInc>
   <timeDataFile name="sin1kHz.dat"/>
</transient>
\end{verbatim}

\noindent
Herein, the different tags have the following meaning:\\
\begin{tabular}{lll} 
$\circ$ & {\bf numSteps}    &  number of time steps\\
$\circ$ & {\bf firstDt }    &  time step size\\
$\circ$ & {\bf stepSaveBeg}  &  number of first step for saving results\\
$\circ$ & {\bf stepSaveEnd}  &  number of last step for saving results\\
$\circ$ & {\bf stepSaveIncr} & increment in steps for saving results\\
$\circ$ & {\bf timeDataFile} & file which holds the varying time values
\end{tabular}
\bigskip

\noindent
Additionally, an input time-file must be provided:
\begin{verbatim} 
<timeDataFile name="sin1kHz.dat"/>
\end{verbatim} 
Therein, you have to define the time values in a column oriented
manner:
\begin{verbatim}
timeVal1          funcVal1
timeVal2          funcVal2
timeVal3          funcVal3
...
\end{verbatim}


\section{List of PDEs}
Inside the config file command
\begin{verbatim}
<pdeList>
   List of PDEs ...
</pdeList>
\end{verbatim}
one has to specify the involved PDEs.  Some possible PDE names are
given in the following table: 

\begin{center}
  \begin{tabular}{l|l}
    {\it Type}      & {\it Description}  \\\hline
    {\bf acoustic}     & Acoustic PDE  \\
    {\bf electrostatic} & Electrostic PDE \\
    {\bf mechanic}     & Mechanics PDE\\    
  \end{tabular}
\end{center}

\noindent
As an example, a coupled electrostatic, mechanical simulation must be
defined as
\begin{verbatim}
<pdeList>
   <mechanic>
      Definitions of Mech-PDE ...
   </mechanic>

   <electrostatic>
      Definitions of Electrostatic-PDE ...
   </electrostatic>
</pdeList>
\end{verbatim}



\section{Geometrical Definitions}
In the top level section {\bf $<$domain$>$} one has to define all
geometrical entities, such as regions and their materials, boundary
nodes, save nodes, ...
\begin{verbatim}
<domain>
  <region name="cantilev" material="silizium"/>
  <region name="gap"      material="air"/>

  <nodes  name="xDof"/>
  <nodes  name="yDof"/>
  <nodes  name="yLoad"/>

  <elements name="mySaveElem"/>

</domain>
\end{verbatim}

The keyword {\bf region} assigns a subdomain called {\bf name} to a material.
Further on, all in the simulation needed nodes (e.g. nodes which hold
a boundary condition) must be defined by the keyword {\bf nodes}.
Needed elements have to be defined by the keyword {\bf elements}.




%%% Local Variables: 
%%% mode: latex
%%% TeX-master: "singlefields"
%%% End: 
